\documentclass[9pt,a4paper]{extarticle}
\usepackage[utf8]{inputenc}
\usepackage[T1]{fontenc}
\usepackage[margin=1.6cm,top=1.4cm,bottom=1.4cm]{geometry}
\usepackage{graphicx,booktabs,array,xcolor,hyperref}
\usepackage{helvet}
\usepackage[indonesian]{babel}
\definecolor{agri}{HTML}{5B7245}
\hypersetup{colorlinks=true,linkcolor=agri,urlcolor=agri}
\setlength{\parskip}{3pt}\setlength{\parindent}{0pt}
\renewcommand{\familydefault}{\sfdefault}
\pagestyle{empty}
\newcommand{\sect}[1]{\vspace{5pt}{\color{agri}\normalsize\bfseries #1}\par\vspace{2pt}}
\newcommand{\tsize}{\footnotesize}
\begin{document}
{\color{agri}\LARGE\bfseries Sinyal Permintaan Pangan Jawa Timur}\\[1pt]
{\normalsize Laporan sinyal permintaan untuk instansi dan lembaga}\\[2pt]
{\footnotesize Periode 1 Agustus sampai 8 September 2026 \,\textbullet\, AgriFlow, \href{https://agriflow.farm}{agriflow.farm}}

\vspace{3pt}\hrule\vspace{4pt}

\sect{Ringkasan}
\begin{itemize}\setlength{\itemsep}{0pt}\setlength{\topsep}{0pt}
\item Dalam 39 hari, \textbf{220 pengunjung} melakukan \textbf{9.530} interaksi dalam \textbf{611 sesi} (rata-rata 3,5 menit per sesi).
\item Perhatian terpusat pada tiga komoditas: cabai rawit, bawang merah, dan beras medium menyerap \textbf{74\%} pemilihan komoditas.
\item Wilayah yang paling banyak dilihat: Kota Surabaya, Banyuwangi, Jember. Untuk cabai rawit, perhatian tertinggi di Malang (25 pengguna-hari), wilayah yang menurut neraca BPS tidak defisit cabai rawit, pola daerah surplus yang mencari pembeli.
\item Skenario yang paling sering diuji: Semeru + banjir (34 kali) dan BBM naik (29 kali) dari 188 simulasi; ini sinyal risiko logistik yang dirasakan pengguna.
\item 335 tampilan prakiraan harga dan 149 unduhan daftar match menunjukkan pemakaian untuk keputusan, bukan sekadar kunjungan.
\end{itemize}

\sect{Apa yang diukur, dan apa yang tidak}
Laporan ini mengukur \textbf{perhatian}: komoditas dan wilayah yang dipilih, prakiraan yang dilihat, skenario yang diuji, dan berkas yang diunduh oleh pengguna dashboard AgriFlow yang telah memberikan persetujuan. Laporan ini tidak mengukur niat beli atau jual, dan tidak memuat data percakapan WhatsApp dalam bentuk apa pun, sesuai WhatsApp Business Solution Terms dan UU PDP No. 27/2022. Tiga sumber digabungkan: interaksi web first-party, neraca surplus-defisit BPS 2022 per kabupaten/kota, dan keluaran engine pencocokan AgriFlow. Setiap pengguna diwakili token harian yang tidak dapat ditautkan lintas hari; sel dengan kurang dari 5 pengguna unik tidak pernah dilaporkan. Satuan ``pengguna-hari'' berarti satu pengguna yang aktif pada satu hari; orang yang sama pada hari berbeda dihitung terpisah, dan itu disengaja.

\begin{minipage}[t]{0.48\linewidth}
\sect{Perhatian per komoditas}
\includegraphics[width=\linewidth]{fig_komoditas.pdf}
\end{minipage}\hfill
\begin{minipage}[t]{0.48\linewidth}
\sect{Pengguna per hari}
\includegraphics[width=\linewidth]{fig_harian.pdf}
\end{minipage}

\begin{minipage}[t]{0.48\linewidth}
{\tsize\begin{tabular}{@{}l r r r r@{}}\toprule
Komoditas & Peng.\,hari & Prakiraan & Simulasi & Unduh \\\midrule
Bawang merah & 380 & 104 & 50 & 40 \\
Cabai rawit & 365 & 95 & 64 & 44 \\
Beras medium & 304 & 82 & 51 & 33 \\
Bawang putih & 117 & 17 & 8 & 12 \\
Beras premium & 108 & 17 & 6 & 7 \\
Cabai merah & 92 & 10 & 5 & 8 \\
Daging ayam & 36 & 5 & 2 & 4 \\
Telur ayam & 25 & 5 & 2 & 1 \\
\bottomrule\end{tabular}}
\end{minipage}\hfill
\begin{minipage}[t]{0.48\linewidth}
\footnotesize Rata-rata 16 pengguna unik per hari; pekan pertama 15, pekan terakhir 20, dengan puncak 43 pengguna pada 2026-09-08. Akhir pekan (abu-abu) sekitar separuh hari kerja. Pada volume ini sel per kabupaten melewati ambang 5 pengguna pada agregasi mingguan, sehingga laporan memakai grain mingguan untuk wilayah dan harian untuk komoditas.

\vspace{3pt}{\tsize\begin{tabular}{@{}l r r@{}}\toprule
Tab & Tampilan & Peng.\,hari \\\midrule
Beranda & 532 & 360 \\
Peta Pasokan & 506 & 356 \\
Rekomendasi Distribusi & 442 & 329 \\
Harga \& Prakiraan & 335 & 268 \\
Simulasi What-if & 265 & 220 \\
Laporan \& KPI & 170 & 152 \\
Bantuan & 90 & 83 \\
Notifikasi & 78 & 74 \\
\bottomrule\end{tabular}}
\end{minipage}

\newpage
\sect{Perhatian per wilayah}
\begin{minipage}[t]{0.5\linewidth}
{\tsize\begin{tabular}{@{}l r r r@{}}\toprule
Kabupaten/Kota & Peng.\,hari & Interaksi & Prakiraan \\\midrule
Kota Surabaya & 54 & 108 & 29 \\
Banyuwangi & 51 & 101 & 26 \\
Jember & 49 & 74 & 16 \\
Malang & 47 & 95 & 26 \\
Gresik & 40 & 79 & 21 \\
Pasuruan & 32 & 53 & 14 \\
Sidoarjo & 31 & 54 & 16 \\
Sampang & 31 & 45 & 12 \\
\bottomrule\end{tabular}}
\end{minipage}\hfill
\begin{minipage}[t]{0.46\linewidth}
{\tsize\begin{tabular}{@{}l l r@{}}\toprule
Komoditas & Wilayah & Peng.\,hari \\\midrule
Bawang merah & Kota Surabaya & 29 \\
Cabai rawit & Malang & 25 \\
Bawang merah & Banyuwangi & 22 \\
Cabai rawit & Kota Surabaya & 22 \\
Cabai rawit & Banyuwangi & 20 \\
Bawang merah & Malang & 18 \\
Cabai rawit & Gresik & 18 \\
Beras medium & Kota Surabaya & 17 \\
\bottomrule\end{tabular}}
\end{minipage}

\sect{Perhatian dan neraca: di mana orang bertanya, dan apakah pasokan ada di sana}
\begin{minipage}[t]{0.5\linewidth}
\vspace{0pt}\includegraphics[width=\linewidth]{fig_scatter.pdf}
\end{minipage}\hfill
\begin{minipage}[t]{0.47\linewidth}
\vspace{0pt}\footnotesize Wilayah dengan perhatian tertinggi adalah Kota Surabaya (54 pengguna-hari, defisit 5.972 ton cabai rawit menurut BPS), Banyuwangi (51, tidak defisit), dan Jember (49, tidak defisit). Kota Surabaya menanggung defisit terbesar (6,0 ribu ton) sekaligus paling banyak dilihat: kuadran defisit tinggi dan perhatian tinggi, prioritas operasi pasar yang paling jelas. Sebaliknya, Sidoarjo (4,3 ribu ton) dan Trenggalek (0,9 ribu ton) defisit besar tetapi hanya dilihat 31 dan 22 pengguna-hari; kuadran itulah yang layak didorong lewat sosialisasi TPID dan operasi pasar. Kuadran perhatian tinggi dan defisit kecil, seperti Banyuwangi dan Jember, menandai daerah surplus yang mencari pembeli, pola yang cocok untuk program penyerapan.

\vspace{4pt}{\tsize\begin{tabular}{@{}l r@{}}\toprule
Skenario yang diuji & Jumlah \\\midrule
Semeru + banjir & 34 \\
BBM naik & 29 \\
Erupsi Semeru & 27 \\
Ramadan & 26 \\
Banjir sentra padi & 24 \\
Ramadan + BBM naik & 24 \\
Suramadu ditutup & 24 \\
\bottomrule\end{tabular}}
\end{minipage}

\sect{Cara memakai laporan ini}
\begin{itemize}\setlength{\itemsep}{0pt}\setlength{\topsep}{0pt}
\item \textbf{Bapanas dan TPID provinsi/kabupaten}: pemantauan mingguan komoditas yang perhatiannya naik sebelum harga bergerak; pemilihan wilayah prioritas operasi pasar dari kuadran defisit tinggi dan perhatian tinggi.
\item \textbf{Kemendag dan Kementan}: bukti permintaan informasi per komoditas untuk menyusun prioritas data terbuka dan jadwal rilis.
\item \textbf{Lembaga keuangan, asuransi, dan pemasok sarana produksi}: peta perhatian per kabupaten sebagai indikator awal aktivitas pasar, digabung dengan neraca BPS.
\item \textbf{Logistik}: skenario gangguan yang paling sering diuji sebagai daftar risiko yang dipersepsikan pelaku.
\end{itemize}

\sect{Batasan}
Sampel awal masih kecil dan berasal dari pengguna yang mengenal AgriFlow, sehingga belum mewakili seluruh pelaku pangan Jawa Timur. Perhatian bukan transaksi. Sel di bawah 5 pengguna disuppress, sehingga wilayah kecil muncul lebih lambat. Riwayat harga yang mendasari prakiraan berhenti pada tanggal pembaruan terakhir sumber Siskaperbapo. Semua batas ini tercantum agar pembaca dapat menimbang angkanya sendiri.

\sect{Berlangganan}
Laporan mingguan atau dua mingguan, format PDF dan CSV, per komoditas atau per wilayah, dengan akses ke tampilan \texttt{demand\_signal\_export} untuk integrasi. Hubungi AgriFlow melalui \href{https://agriflow.farm}{agriflow.farm}.

\vfill{\scriptsize Sumber: interaksi dashboard AgriFlow (persetujuan versi 2026-09-v1), BPS Jawa Timur 2022, engine AgriFlow v1.1. Dokumen ini dihasilkan otomatis dari tampilan agregat; tidak ada data perorangan yang dibaca dalam penyusunannya.}
\end{document}
